\documentclass[a4paper,landscape,8pt]{extarticle}
\usepackage[landscape,margin=0.7cm,top=0.5cm,bottom=0.7cm]{geometry}
\usepackage[utf8]{inputenc}
\usepackage[T1]{fontenc}
\usepackage[ngerman]{babel}
\usepackage{helvet}
\renewcommand{\familydefault}{\sfdefault}
\usepackage{xcolor}
\usepackage{tikz}
\usetikzlibrary{positioning}
\usepackage{enumitem}
\usepackage{multicol}
\usepackage{array}
\usepackage{fancyhdr}
\usepackage{amssymb}

% Farben
\definecolor{calred}{HTML}{E53935}
\definecolor{contactblue}{HTML}{1976D2}
\definecolor{remindergreen}{HTML}{43A047}
\definecolor{lightgray}{HTML}{F5F5F5}
\definecolor{bordergray}{HTML}{BBBBBB}

\pagestyle{fancy}
\fancyhf{}
\renewcommand{\headrulewidth}{0pt}
\fancyfoot[C]{\footnotesize Seite \thepage{} von 3}

\setlength{\parindent}{0pt}
\setlength{\parskip}{2pt}

% Kompakte Listen
\setlist{nosep, leftmargin=1em, itemsep=0pt, topsep=1pt}

\begin{document}

% ============================================================================
% SEITE 1: INFORMATIONSBLATT
% ============================================================================

\noindent{\Large\bfseries Kalender \& Kontakte: Nie mehr vergessen!} \hfill
{\footnotesize\textbf{Lernziele:} Termine eintragen | Erinnerungen | Kontakte anlegen}

\vspace{1mm}
\hrule
\vspace{1mm}

% Drei-Spalten-Layout
\begin{minipage}[t]{0.32\linewidth}
\begin{center}
\colorbox{calred!15}{%
\begin{minipage}{0.95\linewidth}
\vspace{2mm}
\begin{center}
{\Large\bfseries\textcolor{calred}{1. Termin eintragen}}\\[1pt]
{\footnotesize\itshape Nie mehr vergessen}
\end{center}
\vspace{2mm}
\end{minipage}%
}
\end{center}

\vspace{1mm}

\textbf{Kalender-App öffnen:}\\
Buntes Symbol mit Tageszahl

\textbf{So trägst du einen Termin ein:}
\begin{enumerate}
\item Kalender-App öffnen
\item \textbf{+} oben rechts antippen
\item \textbf{Titel} eingeben (z.B. ``Arzt Dr. Müller'')
\item ``Beginnt'' antippen $\rightarrow$ \textbf{Datum und Uhrzeit} wählen
\item ``Hinweis'' antippen $\rightarrow$ z.B. ``1 Tag davor''
\item ``Hinzufügen'' oder ``Fertig''
\end{enumerate}

\textbf{Ansichten:}
\begin{itemize}
\item \textbf{Tag:} Ein Tag im Detail
\item \textbf{Monat:} Übersicht mit Punkten
\item Unten umschalten
\end{itemize}

\textbf{Praktische Termine:}
\begin{itemize}
\item Arzttermine (mit Adresse!)
\item Geburtstage (jährlich)
\item Müllabfuhr
\end{itemize}

\end{minipage}%
\hfill%
\begin{minipage}[t]{0.32\linewidth}
\begin{center}
\colorbox{remindergreen!15}{%
\begin{minipage}{0.95\linewidth}
\vspace{2mm}
\begin{center}
{\Large\bfseries\textcolor{remindergreen}{2. Erinnerungen}}\\[1pt]
{\footnotesize\itshape Das Handy erinnert dich}
\end{center}
\vspace{2mm}
\end{minipage}%
}
\end{center}

\vspace{1mm}

\textbf{Warum Erinnerungen?}\\
Das Handy \textbf{stupst dich an}, bevor der Termin beginnt.

\textbf{Erinnerung einstellen:}
\begin{enumerate}
\item Im Termin: ``Hinweis'' antippen
\item Zeitpunkt wählen:
\begin{itemize}
\item ``Zur Startzeit''
\item ``5 Minuten davor''
\item ``1 Stunde davor''
\item ``\textbf{1 Tag davor}'' (für Arzt!)
\item ``1 Woche davor''
\end{itemize}
\end{enumerate}

\textbf{Was passiert dann?}
\begin{itemize}
\item Handy zeigt \textbf{Benachrichtigung}
\item Ton oder Vibration
\item Auch auf dem Sperrbildschirm
\end{itemize}

\vspace{2mm}
\fcolorbox{remindergreen}{remindergreen!5}{%
\begin{minipage}{0.92\linewidth}
\footnotesize
\textbf{Tipp:} Für Arzttermine ``1 Tag davor'' einstellen, dann kannst du dich vorbereiten!
\end{minipage}%
}

\end{minipage}%
\hfill%
\begin{minipage}[t]{0.32\linewidth}
\begin{center}
\colorbox{contactblue!15}{%
\begin{minipage}{0.95\linewidth}
\vspace{2mm}
\begin{center}
{\Large\bfseries\textcolor{contactblue}{3. Kontakte anlegen}}\\[1pt]
{\footnotesize\itshape Dein digitales Adressbuch}
\end{center}
\vspace{2mm}
\end{minipage}%
}
\end{center}

\vspace{1mm}

\textbf{Kontakte-App:}\\
Symbol mit Silhouette (Telefon $\rightarrow$ Kontakte)

\textbf{Neuen Kontakt anlegen:}
\begin{enumerate}
\item Kontakte-App öffnen
\item \textbf{+} oben rechts
\item Vorname, Nachname eingeben
\item ``Telefon hinzufügen''
\item Optional: E-Mail, Adresse, Geburtstag
\item ``Fertig''
\end{enumerate}

\textbf{Kontakt nutzen:}
\begin{itemize}
\item Antippen $\rightarrow$ \textbf{Anrufen}
\item Nachricht-Symbol $\rightarrow$ \textbf{SMS/WhatsApp}
\item Mail-Symbol $\rightarrow$ \textbf{E-Mail}
\end{itemize}

\textbf{Favoriten:}
\begin{itemize}
\item Kontakt öffnen
\item ``Zu Favoriten hinzufügen''
\item Erscheinen ganz oben in der Telefon-App
\end{itemize}

\end{minipage}

\vspace{1mm}

% Zusammenfassung
\begin{center}
\fcolorbox{bordergray}{lightgray}{%
\begin{minipage}{0.98\linewidth}
\centering
\vspace{1.5mm}
\textbf{Zusammenfassung:}
\textcolor{calred}{\textbf{Termin}} = + Symbol, Titel, Datum, Uhrzeit ~$|$~
\textcolor{remindergreen}{\textbf{Erinnerung}} = ``Hinweis'' einstellen ~$|$~
\textcolor{contactblue}{\textbf{Kontakt}} = + Symbol, Name, Telefonnummer
\vspace{1.5mm}
\end{minipage}%
}
\end{center}

\newpage

% ============================================================================
% SEITE 2: ÜBUNGEN
% ============================================================================

\begin{center}
{\LARGE\bfseries Übungen: Kalender \& Kontakte}\\[3pt]
{\normalsize Schau dir Seite 1 an und beantworte die Fragen}
\end{center}

\vspace{3mm}
\hrule
\vspace{4mm}

\begin{multicols}{2}

\textbf{\large Teil A: Grundlagen}

Richtig oder Falsch? Kreuze an.

\vspace{2mm}

\renewcommand{\arraystretch}{1.6}
\begin{tabular}{@{}p{8cm}cc@{}}
& R & F \\
1. Termine trägt man mit dem + Symbol ein. & $\square$ & $\square$ \\
2. Erinnerungen können nur 5 Minuten vorher kommen. & $\square$ & $\square$ \\
3. Kontakte kann man zu Favoriten machen. & $\square$ & $\square$ \\
4. Aus einem Kontakt kann man direkt anrufen. & $\square$ & $\square$ \\
5. Der Kalender kann Geburtstage speichern. & $\square$ & $\square$ \\
\end{tabular}
\renewcommand{\arraystretch}{1}

\vspace{6mm}

\textbf{\large Teil B: Termin eintragen}

Nummeriere die Schritte (1--5):

\vspace{2mm}

\renewcommand{\arraystretch}{1.8}
\begin{tabular}{@{}cp{7cm}@{}}
$\square$ & ``Fertig'' antippen \\
$\square$ & + Symbol antippen \\
$\square$ & Erinnerung (``Hinweis'') einstellen \\
$\square$ & Kalender-App öffnen \\
$\square$ & Titel, Datum und Uhrzeit eingeben \\
\end{tabular}
\renewcommand{\arraystretch}{1}

\vspace{6mm}

\textbf{\large Teil C: Erinnerungszeiten}

Welche Erinnerung passt? Verbinde:

\vspace{2mm}

\renewcommand{\arraystretch}{1.8}
\begin{tabular}{@{}p{4cm}p{4cm}@{}}
Arzttermin & $\circ$ ~~~~ $\circ$ 5 Minuten vorher \\
Telefonkonferenz & $\circ$ ~~~~ $\circ$ 1 Tag vorher \\
Geburtstag (Geschenk!) & $\circ$ ~~~~ $\circ$ 1 Woche vorher \\
\end{tabular}
\renewcommand{\arraystretch}{1}

\columnbreak

\textbf{\large Teil D: Praktisch üben}

Führe diese Aufgaben auf deinem Gerät durch:

\vspace{3mm}

\textbf{Aufgabe 1: Termin eintragen}\\
Trage einen Termin ein: ``Übungs-Termin'' für morgen um 10:00 Uhr mit Erinnerung ``1 Stunde vorher''.

\vspace{2mm}
$\square$ Erledigt

\vspace{5mm}

\textbf{Aufgabe 2: Kontakt anlegen}\\
Lege einen neuen Kontakt an (z.B. einen Bekannten oder fiktiv ``Test Person'').

\vspace{2mm}
$\square$ Erledigt

\vspace{5mm}

\textbf{Aufgabe 3: Favorit hinzufügen}\\
Füge einen wichtigen Kontakt zu deinen Favoriten hinzu.

\vspace{2mm}
$\square$ Erledigt

\vspace{5mm}

\textbf{Aufgabe 4: Aus Kontakt anrufen}\\
Öffne einen Kontakt und finde die Anrufen-Funktion (nicht antippen, nur finden).

\vspace{2mm}
$\square$ Gefunden

\vspace{5mm}

\textbf{Bonus: Geburtstag eintragen}\\
Trage einen wichtigen Geburtstag in den Kalender ein (jährlich wiederholen).

\vspace{2mm}
$\square$ Erledigt

\end{multicols}

\newpage

% ============================================================================
% SEITE 3: CHECKLISTE
% ============================================================================

\begin{center}
{\LARGE\bfseries Checkliste: Das habe ich verstanden}\\[3pt]
{\normalsize Geh die Liste durch und hake ab, was du sicher weißt}
\end{center}

\vspace{3mm}
\hrule
\vspace{6mm}

% Drei Boxen nebeneinander
\begin{minipage}[t]{0.31\linewidth}
\begin{center}
\fcolorbox{calred}{calred!10}{%
\begin{minipage}{0.92\linewidth}
\vspace{3mm}
\begin{center}
{\large\bfseries\textcolor{calred}{Kalender}}
\end{center}
\vspace{2mm}
\begin{itemize}[label=$\square$, itemsep=5pt]
\item Ich kann die \textbf{Kalender-App} öffnen
\item Ich kann einen \textbf{Termin eintragen}
\item Ich weiß, wie ich \textbf{Datum und Uhrzeit} setze
\item Ich kann zwischen \textbf{Ansichten} wechseln
\end{itemize}
\vspace{3mm}
\end{minipage}%
}
\end{center}
\end{minipage}%
\hfill%
\begin{minipage}[t]{0.31\linewidth}
\begin{center}
\fcolorbox{remindergreen}{remindergreen!10}{%
\begin{minipage}{0.92\linewidth}
\vspace{3mm}
\begin{center}
{\large\bfseries\textcolor{remindergreen}{Erinnerungen}}
\end{center}
\vspace{2mm}
\begin{itemize}[label=$\square$, itemsep=5pt]
\item Ich weiß, wo ich \textbf{``Hinweis''} finde
\item Ich kann die \textbf{Erinnerungszeit} wählen
\item Ich verstehe die \textbf{verschiedenen Zeiten}
\item Ich weiß: Das Handy \textbf{erinnert mich}
\end{itemize}
\vspace{3mm}
\end{minipage}%
}
\end{center}
\end{minipage}%
\hfill%
\begin{minipage}[t]{0.31\linewidth}
\begin{center}
\fcolorbox{contactblue}{contactblue!10}{%
\begin{minipage}{0.92\linewidth}
\vspace{3mm}
\begin{center}
{\large\bfseries\textcolor{contactblue}{Kontakte}}
\end{center}
\vspace{2mm}
\begin{itemize}[label=$\square$, itemsep=5pt]
\item Ich kann einen \textbf{Kontakt anlegen}
\item Ich kann aus einem Kontakt \textbf{anrufen}
\item Ich kann \textbf{Favoriten} hinzufügen
\item Kontakte sind mein \textbf{digitales Adressbuch}
\end{itemize}
\vspace{3mm}
\end{minipage}%
}
\end{center}
\end{minipage}

\vspace{8mm}

% Gesamtverständnis
\begin{center}
\fcolorbox{black}{lightgray}{%
\begin{minipage}{0.95\linewidth}
\vspace{4mm}
\begin{center}
{\large\bfseries Gesamtverständnis}
\end{center}
\vspace{3mm}
\begin{multicols}{2}
\begin{itemize}[label=$\square$, itemsep=6pt]
\item Ich kann \textbf{Termine} mit Erinnerung eintragen
\item Ich vergesse keine \textbf{wichtigen Termine} mehr
\item Ich kann \textbf{Kontakte} anlegen und nutzen
\item Ich kann direkt aus einem Kontakt \textbf{anrufen}
\item Ich fühle mich \textbf{organisierter}
\end{itemize}
\end{multicols}
\vspace{3mm}
\end{minipage}%
}
\end{center}

\vspace{8mm}

% Notfall-Notizen
\begin{center}
\fcolorbox{bordergray}{white}{%
\begin{minipage}{0.95\linewidth}
\vspace{4mm}
\begin{center}
{\large\bfseries Meine wichtigen Termine \& Kontakte}\\[3pt]
{\footnotesize (Trage hier ein, was du digital speichern möchtest)}
\end{center}
\vspace{4mm}
\renewcommand{\arraystretch}{2}
\begin{tabular}{@{}p{0.45\linewidth}p{0.45\linewidth}@{}}
\textbf{Wichtige regelmäßige Termine:} & \textbf{Wichtige Kontakte als Favorit:} \\
1. \dotfill & 1. \dotfill \\[2mm]
2. \dotfill & 2. \dotfill \\[2mm]
3. \dotfill & 3. \dotfill \\[2mm]
4. \dotfill & 4. \dotfill \\
\end{tabular}
\renewcommand{\arraystretch}{1}
\vspace{4mm}
\end{minipage}%
}
\end{center}

\vfill

\begin{center}
\footnotesize\textit{Stand: \today{} -- Kalender + Erinnerungen = Nie mehr Termine vergessen!}
\end{center}

\end{document}
